% =====================================================================
%  papers/shared/preamble.tex
%  Shared preamble for paper1 and paper2.  Owned by T-005 (scaffold).
%
%  Package set is restricted to the verified-working list:
%    amsmath, amssymb, amsthm, booktabs, graphicx, hyperref, natbib, geometry
%  Do not add packages without compiling to confirm they resolve.
%
%  AUTHORING RULE (D-005): this file and every other .tex in this tree are
%  written with the Write/Edit tools.  Bash heredocs collapse \\ to \ in this
%  environment and silently corrupt tabulars.
% =====================================================================

\usepackage[a4paper,margin=1.1in]{geometry}

\usepackage{amsmath}
\usepackage{amssymb}
\usepackage{amsthm}
\usepackage{booktabs}
\usepackage{graphicx}
\usepackage[authoryear,round]{natbib}

% hyperref last (except for its own configuration below).
\usepackage[
  colorlinks=true,
  linkcolor=black,
  citecolor=black,
  urlcolor=black,
  bookmarksnumbered=true,
  bookmarksopen=true,
  pdfstartview=FitH
]{hyperref}

% Figures live next to each paper's main .tex, so a bare relative path works
% for both documents when latexmk runs inside the paper directory.
\graphicspath{{figures/}{./figures/}}

\bibliographystyle{plainnat}
\setcitestyle{authoryear,round}

% ---------------------------------------------------------------------
%  Theorem-like environments
% ---------------------------------------------------------------------
\theoremstyle{plain}
\newtheorem{theorem}{Theorem}[section]
\newtheorem{proposition}[theorem]{Proposition}
\newtheorem{lemma}[theorem]{Lemma}
\newtheorem{corollary}[theorem]{Corollary}

\theoremstyle{definition}
\newtheorem{definition}[theorem]{Definition}
\newtheorem{example}[theorem]{Example}

\theoremstyle{remark}
\newtheorem{remark}[theorem]{Remark}

% ---------------------------------------------------------------------
%  \begin{keyinsight} ... \end{keyinsight}
%
%  Marks a contribution or a load-bearing idea so a reader skimming for the
%  new result can find it.  Built from base LaTeX primitives only -- no
%  framing package is in the verified set.
%  Optional argument overrides the label, e.g. \begin{keyinsight}[What is new]
% ---------------------------------------------------------------------
\newenvironment{keyinsight}[1][Key insight]{%
  \par\medskip
  \begin{list}{}{%
    \setlength{\leftmargin}{1.5em}%
    \setlength{\rightmargin}{1.5em}%
    \setlength{\listparindent}{0pt}%
    \setlength{\parsep}{0pt}%
    \setlength{\topsep}{0pt}%
    \setlength{\partopsep}{0pt}%
  }%
  \item[]%
  \noindent\rule{\linewidth}{0.9pt}\par\smallskip
  \noindent\textbf{#1.}\enspace\ignorespaces
}{%
  \par\smallskip\noindent\rule{\linewidth}{0.9pt}%
  \end{list}%
  \medskip
}

% ---------------------------------------------------------------------
%  Labelled section scheme (AC-012)
%
%  Every top-level section carries a role tag that survives into the table
%  of contents and the PDF bookmarks, so a reader can locate the new result
%  without reading linearly.  Roles:
%    exposition -- background a reader may already know and may skip
%    evidence   -- computation and measurement supporting the claim
%    contribution -- what this paper adds
%    limitations  -- what the claim does not cover
% ---------------------------------------------------------------------
\newcommand{\roletag}[1]{%
  \texorpdfstring{\ {\normalfont\small\textsc{[#1]}}}{ [#1]}%
}
\newcommand{\EXPOSITION}{\roletag{exposition}}
\newcommand{\EVIDENCE}{\roletag{evidence}}
\newcommand{\CONTRIBUTION}{\roletag{contribution}}
\newcommand{\LIMITATIONS}{\roletag{limitations}}

% ---------------------------------------------------------------------
%  Glossary machinery (AC-003)
%
%  \defterm{key}{printed term}
%      Marks the defining occurrence of a term of art in the running text.
%      It sets the term in bold italic and plants a label.
%
%  \gentry{key}{Headword}{plain-language definition}
%      A glossary line.  It back-references the section and page where
%      \defterm{key}{...} was placed.
%
%  INVARIANT: every \gentry key should have a matching \defterm in a file
%  that the paper actually \inputs.  If it does not, \gentry degrades to
%  "Definition pending" and logs a LaTeX warning rather than emitting ??,
%  so a missing anchor cannot silently break AC-001 for the other paper.
%  Grep the build log for "glossary anchor" to find any that are pending.
% ---------------------------------------------------------------------
\newcommand{\defterm}[2]{\textbf{\textit{#2}}\label{term:#1}}

\makeatletter
\newcommand{\gentry}[3]{%
  \item[#2]\ #3\ %
  \@ifundefined{r@term:#1}{%
    \@latex@warning{Missing glossary anchor for term '#1'%
      \MessageBreak (no \string\defterm{#1}{...} in this document)}%
    \emph{(Definition pending.)}%
  }{%
    \emph{(Defined in Section~\ref{term:#1}, page~\pageref{term:#1}.)}%
  }%
}
\makeatother

% ---------------------------------------------------------------------
%  Small conveniences used across both papers
% ---------------------------------------------------------------------
\newcommand{\Wmat}{\mathbf{W}}
\newcommand{\Gistar}{G_i^{*}}
\newcommand{\MoranI}{I}
\newcommand{\Ex}{\mathbb{E}}
\newcommand{\Var}{\operatorname{Var}}

% Author metadata shared by both papers.
\newcommand{\paperauthor}{Matthew Scott Tan}
